\documentclass[11pt]{article}
% The preamble every project's paper shares: % The preamble every project's paper shares: % The preamble every project's paper shares: \input{../../preamble} from
% projects/<project>/paper/. Boards come from chessboard and xskak, plots from
% pgfplots; every number a paper prints comes from its generated.tex.
\usepackage[LSBC4,T1]{fontenc}
\usepackage[utf8]{inputenc}
\usepackage{lmodern}
\usepackage[margin=1in]{geometry}
\usepackage{microtype}
\usepackage{amsmath}
\usepackage{booktabs}
\usepackage{tabularx}
\usepackage{array}
\usepackage{xcolor}
\usepackage[obeyspaces,spaces]{url}
\usepackage{xskak}
\usepackage{pgfplots}
\pgfplotsset{compat=1.18}
\usepackage{caption}
\usepackage{subcaption}
\usepackage{enumitem}
\usepackage[section]{placeins}
\usepackage[hidelinks]{hyperref}

\setchessboard{boardfontencoding=LSBC4, boardfontfamily=skaknew, blackfieldcolor=black!22,
  whitefieldcolor=white, setfontcolors, showmover=false, label=true, labelfontsize=6pt,
  boardfontsize=14pt}

% A move or line in UCI, and a position in FEN, both breakable at a slash.
\newcommand{\uci}[1]{\texttt{#1}}
\newcommand{\fen}[1]{{\footnotesize\path{#1}}}
\newcolumntype{L}{>{\raggedright\arraybackslash}X}

\setlength{\parskip}{0.4em}
\setlength{\parindent}{0pt}
 from
% projects/<project>/paper/. Boards come from chessboard and xskak, plots from
% pgfplots; every number a paper prints comes from its generated.tex.
\usepackage[LSBC4,T1]{fontenc}
\usepackage[utf8]{inputenc}
\usepackage{lmodern}
\usepackage[margin=1in]{geometry}
\usepackage{microtype}
\usepackage{amsmath}
\usepackage{booktabs}
\usepackage{tabularx}
\usepackage{array}
\usepackage{xcolor}
\usepackage[obeyspaces,spaces]{url}
\usepackage{xskak}
\usepackage{pgfplots}
\pgfplotsset{compat=1.18}
\usepackage{caption}
\usepackage{subcaption}
\usepackage{enumitem}
\usepackage[section]{placeins}
\usepackage[hidelinks]{hyperref}

\setchessboard{boardfontencoding=LSBC4, boardfontfamily=skaknew, blackfieldcolor=black!22,
  whitefieldcolor=white, setfontcolors, showmover=false, label=true, labelfontsize=6pt,
  boardfontsize=14pt}

% A move or line in UCI, and a position in FEN, both breakable at a slash.
\newcommand{\uci}[1]{\texttt{#1}}
\newcommand{\fen}[1]{{\footnotesize\path{#1}}}
\newcolumntype{L}{>{\raggedright\arraybackslash}X}

\setlength{\parskip}{0.4em}
\setlength{\parindent}{0pt}
 from
% projects/<project>/paper/. Boards come from chessboard and xskak, plots from
% pgfplots; every number a paper prints comes from its generated.tex.
\usepackage[LSBC4,T1]{fontenc}
\usepackage[utf8]{inputenc}
\usepackage{lmodern}
\usepackage[margin=1in]{geometry}
\usepackage{microtype}
\usepackage{amsmath}
\usepackage{booktabs}
\usepackage{tabularx}
\usepackage{array}
\usepackage{xcolor}
\usepackage[obeyspaces,spaces]{url}
\usepackage{xskak}
\usepackage{pgfplots}
\pgfplotsset{compat=1.18}
\usepackage{caption}
\usepackage{subcaption}
\usepackage{enumitem}
\usepackage[section]{placeins}
\usepackage[hidelinks]{hyperref}

\setchessboard{boardfontencoding=LSBC4, boardfontfamily=skaknew, blackfieldcolor=black!22,
  whitefieldcolor=white, setfontcolors, showmover=false, label=true, labelfontsize=6pt,
  boardfontsize=14pt}

% A move or line in UCI, and a position in FEN, both breakable at a slash.
\newcommand{\uci}[1]{\texttt{#1}}
\newcommand{\fen}[1]{{\footnotesize\path{#1}}}
\newcolumntype{L}{>{\raggedright\arraybackslash}X}

\setlength{\parskip}{0.4em}
\setlength{\parindent}{0pt}

% Written by projects/period-video-coaching/paper/generate.mjs; do not edit by hand.

\newcommand{\CandidateBullet}{%
\begin{quote}\small
The video coaches the way a strong coach works with a player at his level, in every board scene that teaches a move of his, from the evidence alone and inside the period video grammar's chapters and rules:
\begin{itemize}[nosep]
\item ask before you tell: the scene's first pause asks what he would play there, open (``What would you play here? Pick two moves.''), before any question that narrows it to a kind of move and before any arrow;
\item name the habit, then its replacement: say what kind of move he kept playing and how (its scan step, the seconds it got, what it left behind, as the dailies and the archive say) and how often it came back this period, and show at the board how the rule that answers it replaces that habit;
\item and then what: after the better move, play or say their best reply and why the move still holds, so the line stops where its gain is clear (the turning point's engine line in coach-log/review-analysis/<date>.json, or \texttt{npm run -{}-silent coach -{}- engine best <fen>});
\item pitch it at his level: say whether this is a miss most players at his rating make or one he could have found, from how often players at his band play the best move there (the \texttt{find} the day's Maia step archived on that turning point in coach-log/review-analysis/<date>.json, or \texttt{npm run -{}-silent coach -{}- engine find <fen> <uci> -{}-rapid <his rating>}), in words and never as a percentage (under about 15\%: a pattern to learn; half or more: his to find); on a hung piece, the mirror instead, as the video grammar says;
\item praise with a board: on a move he got right, name what it saw.
\end{itemize}
\end{quote}
}
\newcommand{\CaseSpan}{weeks 35 to 42 and the months of August and September 2026}
\newcommand{\AllPeriods}{ten}
\newcommand{\Cases}{ten}
\newcommand{\CasesCap}{Ten}
\newcommand{\WeeksN}{eight}
\newcommand{\MonthsN}{two}
\newcommand{\AppCommit}{3eb8afd0c1}
\newcommand{\RefFailedN}{nine}
\newcommand{\RefPassed}{week 42}
\newcommand{\RefErrorsMin}{22}
\newcommand{\RefErrorsMax}{271}
\newcommand{\TrainN}{five}
\newcommand{\HeldN}{five}
\newcommand{\TrainCases}{weeks 35, 37, 39 and 41, and August}
\newcommand{\HeldCases}{weeks 36, 38, 40 and 42, and September}
\newcommand{\WeekGamesMin}{19}
\newcommand{\WeekGamesMax}{59}
\newcommand{\MonthGames}{307 (August) and 172 (September)}
\newcommand{\WriterRuns}{24}
\newcommand{\NullN}{four}
\newcommand{\NullCases}{weeks 35, 36, 37 and 38}
\newcommand{\WordsPerSecond}{2.6}
\newcommand{\MinDecoyWords}{40}
\newcommand{\Spent}{225}
\newcommand{\Cap}{276}
\newcommand{\PilotCap}{50}
\newcommand{\WriterSpend}{155}
\newcommand{\GraderSpend}{63}
\newcommand{\LostSpend}{6.50}
\newcommand{\SevenDayFirst}{55}
\newcommand{\SevenDayLast}{68}
\newcommand{\PublishedN}{ten}
\newcommand{\PubOpenQuestion}{15}
\newcommand{\PubHabitNamed}{90}
\newcommand{\PubReplacement}{81}
\newcommand{\PubThenWhat}{19}
\newcommand{\PubLevel}{38}
\newcommand{\PubLevelTrue}{63}
\newcommand{\PubPraiseNamed}{0}
\newcommand{\PubOpenQuestionAll}{30}
\newcommand{\PubMoments}{2.1}
\newcommand{\PubScenes}{3.4}
\newcommand{\BaseOpenQuestion}{16}
\newcommand{\BaseHabitNamed}{84}
\newcommand{\BaseReplacement}{79}
\newcommand{\BaseThenWhat}{28}
\newcommand{\BaseLevel}{11}
\newcommand{\BaseLevelTrue}{100}
\newcommand{\BasePraiseNamed}{0}
\newcommand{\BaseOpenQuestionAll}{38}
\newcommand{\BaseMoments}{1.9}
\newcommand{\BaseScenes}{3.6}
\newcommand{\CandOpenQuestion}{95}
\newcommand{\CandHabitNamed}{76}
\newcommand{\CandReplacement}{74}
\newcommand{\CandThenWhat}{83}
\newcommand{\CandLevel}{65}
\newcommand{\CandLevelTrue}{90}
\newcommand{\CandPraiseNamed}{100}
\newcommand{\CandOpenQuestionAll}{98}
\newcommand{\CandMoments}{3.1}
\newcommand{\CandScenes}{6.0}
\newcommand{\PubSplit}{0.0}
\newcommand{\PubPositions}{2.1}
\newcommand{\PubToldHabitNamed}{90}
\newcommand{\PubToldReplacement}{81}
\newcommand{\PubToldLevel}{38}
\newcommand{\BaseSplit}{0.0}
\newcommand{\BasePositions}{1.9}
\newcommand{\BaseToldHabitNamed}{84}
\newcommand{\BaseToldReplacement}{79}
\newcommand{\BaseToldLevel}{11}
\newcommand{\CandSplit}{1.2}
\newcommand{\CandPositions}{2.1}
\newcommand{\CandToldHabitNamed}{100}
\newcommand{\CandToldReplacement}{100}
\newcommand{\CandToldLevel}{95}
\newcommand{\MovesTable}{%
\begin{tabular}{lrrr}
\toprule
In the moment scenes & Published & Today's prompts & Candidate \\
\midrule
Asks what he would play before any hint & 15\% {\footnotesize (20)} & 16\% {\footnotesize (19)} & 95\% {\footnotesize (21)} \\
Names the habit behind the mistake & 90\% {\footnotesize (21)} & 84\% {\footnotesize (19)} & 76\% {\footnotesize (29)} \\
Ends on the habit that replaces it & 81\% {\footnotesize (21)} & 79\% {\footnotesize (19)} & 74\% {\footnotesize (31)} \\
Plays out their best reply to the better move & 19\% {\footnotesize (21)} & 28\% {\footnotesize (18)} & 83\% {\footnotesize (24)} \\
Says whether the miss is normal at his level & 38\% {\footnotesize (21)} & 11\% {\footnotesize (19)} & 65\% {\footnotesize (31)} \\
\quad and that agrees with the evidence & 63\% {\footnotesize (8)} & 100\% {\footnotesize (2)} & 90\% {\footnotesize (20)} \\
\midrule
\multicolumn{4}{l}{\emph{Counted once per position taught}} \\
Names the habit behind the mistake & 90\% {\footnotesize (21)} & 84\% {\footnotesize (19)} & 100\% {\footnotesize (21)} \\
Ends on the habit that replaces it & 81\% {\footnotesize (21)} & 79\% {\footnotesize (19)} & 100\% {\footnotesize (21)} \\
Says whether the miss is normal at his level & 38\% {\footnotesize (21)} & 11\% {\footnotesize (19)} & 95\% {\footnotesize (21)} \\
\midrule
Praise names what the good move saw & -- & -- & 100\% {\footnotesize (11)} \\
Any board scene: an open question first & 30\% {\footnotesize (33)} & 38\% {\footnotesize (34)} & 98\% {\footnotesize (46)} \\
\bottomrule
\end{tabular}
}
\newcommand{\ShareBase}{44}
\newcommand{\ShareCand}{79}
\newcommand{\ShareDiff}{+0.36 {\footnotesize (+0.24 to +0.47)}}
\newcommand{\BaseHeard}{6.4}
\newcommand{\BaseWords}{873}
\newcommand{\BaseCost}{6.13}
\newcommand{\BaseMinutes}{18}
\newcommand{\BaseRounds}{6.3}
\newcommand{\BasePassed}{10}
\newcommand{\BaseRuns}{10}
\newcommand{\BaseBoardScenes}{3.6}
\newcommand{\BaseQuestions}{4.2}
\newcommand{\BaseAllScenes}{11.4}
\newcommand{\BaseWeekHeard}{6.3}
\newcommand{\BaseWeekUnder}{zero}
\newcommand{\BaseWeekIn}{eight}
\newcommand{\BaseWeekOver}{zero}
\newcommand{\BaseMonthHeard}{6.5}
\newcommand{\BaseMonthUnder}{one}
\newcommand{\BaseMonthIn}{one}
\newcommand{\BaseMonthOver}{zero}
\newcommand{\CandHeard}{8.6}
\newcommand{\CandWords}{1165}
\newcommand{\CandCost}{6.48}
\newcommand{\CandMinutes}{19}
\newcommand{\CandRounds}{7.1}
\newcommand{\CandPassed}{10}
\newcommand{\CandRuns}{10}
\newcommand{\CandBoardScenes}{6.0}
\newcommand{\CandQuestions}{9.7}
\newcommand{\CandAllScenes}{13.5}
\newcommand{\CandWeekHeard}{8.9}
\newcommand{\CandWeekUnder}{zero}
\newcommand{\CandWeekIn}{one}
\newcommand{\CandWeekOver}{seven}
\newcommand{\CandMonthHeard}{7.5}
\newcommand{\CandMonthUnder}{one}
\newcommand{\CandMonthIn}{zero}
\newcommand{\CandMonthOver}{one}
\newcommand{\PubHeard}{7.6}
\newcommand{\PubWeekHeard}{7.9}
\newcommand{\WeekTargetLo}{five}
\newcommand{\WeekTargetHi}{eight}
\newcommand{\PubMonthHeard}{6.6}
\newcommand{\MonthTargetLo}{six}
\newcommand{\MonthTargetHi}{nine}
\newcommand{\BaseRendered}{5.1}
\newcommand{\BaseWeekRendered}{5.1}
\newcommand{\BaseWeekRenderedUnder}{four}
\newcommand{\BaseWeekRenderedIn}{four}
\newcommand{\BaseWeekRenderedOver}{zero}
\newcommand{\BaseMonthRendered}{5.1}
\newcommand{\BaseMonthRenderedUnder}{two}
\newcommand{\BaseMonthRenderedIn}{zero}
\newcommand{\BaseMonthRenderedOver}{zero}
\newcommand{\CandRendered}{6.8}
\newcommand{\CandWeekRendered}{7.0}
\newcommand{\CandWeekRenderedUnder}{zero}
\newcommand{\CandWeekRenderedIn}{seven}
\newcommand{\CandWeekRenderedOver}{one}
\newcommand{\CandMonthRendered}{5.8}
\newcommand{\CandMonthRenderedUnder}{one}
\newcommand{\CandMonthRenderedIn}{one}
\newcommand{\CandMonthRenderedOver}{zero}
\newcommand{\RenderedShare}{79}
\newcommand{\EstimateOff}{6}
\newcommand{\RenderedRatio}{1.35}
\newcommand{\Voice}{af\_heart}
\newcommand{\VoiceSpeed}{0.97}
\newcommand{\LengthTable}{%
\begin{tabular}{lcrrrrrrr}
\toprule
 & Grammar & Published & \multicolumn{3}{c}{Today's prompts} & \multicolumn{3}{c}{Candidate} \\
\cmidrule(lr){4-6}\cmidrule(lr){7-9}
 & minutes & heard once & heard once & rendered & inside & heard once & rendered & inside \\
\midrule
Weeks & 5 to 8 & 7.9 & 6.3 & 5.1 & 4 of 8 & 8.9 & 7.0 & 7 of 8 \\
Months & 6 to 9 & 6.6 & 6.5 & 5.1 & 0 of 2 & 7.5 & 5.8 & 1 of 2 \\
\bottomrule
\end{tabular}
}
\newcommand{\MonthLengthText}{rendered, the candidate's August and September videos last 4.4 and 7.3 minutes and today's 5.8 and 4.4, against the grammar's six to nine}
\newcommand{\RecentWeeks}{weeks 39, 40, 41 and 42}
\newcommand{\PubRecentHeard}{6.5}
\newcommand{\BaseRecentHeard}{6.6}
\newcommand{\LengthCap}{1.3}
\newcommand{\LengthRatio}{1.37}
\newcommand{\LengthRatioInterval}{1.37 {\footnotesize (1.20 to 1.52)}}
\newcommand{\TrainTruth}{+0.20 {\footnotesize ($-$0.40 to +0.80)}}
\newcommand{\TrainCoaching}{+1.00 {\footnotesize (+1.00 to +1.00)}}
\newcommand{\TrainListening}{$-$0.20 {\footnotesize ($-$0.80 to +0.40)}}
\newcommand{\TrainOverall}{+0.60 {\footnotesize ($-$0.20 to +1.00)}}
\newcommand{\HeldTruth}{$-$0.20 {\footnotesize ($-$0.80 to +0.40)}}
\newcommand{\HeldCoaching}{+1.00 {\footnotesize (+1.00 to +1.00)}}
\newcommand{\HeldListening}{0.00 {\footnotesize ($-$0.80 to +0.80)}}
\newcommand{\HeldOverall}{+0.60 {\footnotesize (+0.20 to +1.00)}}
\newcommand{\AllTruth}{0.00 {\footnotesize ($-$0.50 to +0.50)}}
\newcommand{\AllCoaching}{+1.00 {\footnotesize (+1.00 to +1.00)}}
\newcommand{\AllListening}{$-$0.10 {\footnotesize ($-$0.60 to +0.40)}}
\newcommand{\AllOverall}{+0.60 {\footnotesize (+0.20 to +1.00)}}
\newcommand{\AllNullCoaching}{0.00 {\footnotesize ($-$0.75 to +0.75)}}
\newcommand{\TruthWins}{three}
\newcommand{\TruthLosses}{three}
\newcommand{\TruthTies}{four}
\newcommand{\CoachingWins}{ten}
\newcommand{\CoachingLosses}{zero}
\newcommand{\CoachingTies}{zero}
\newcommand{\ListeningWins}{three}
\newcommand{\ListeningLosses}{four}
\newcommand{\ListeningTies}{three}
\newcommand{\OverallWins}{seven}
\newcommand{\OverallLosses}{one}
\newcommand{\OverallTies}{two}
\newcommand{\NullDecisive}{two}
\newcommand{\Decided}{102}
\newcommand{\SecondShare}{54}
\newcommand{\SameN}{ten}
\newcommand{\SameBroken}{zero}
\newcommand{\DecoyN}{ten}
\newcommand{\DecoyCoaching}{zero}
\newcommand{\DecoyLostListening}{ten}
\newcommand{\DecoyLostOverall}{ten}
\newcommand{\DecoyLostCoaching}{nine}
\newcommand{\DecoyAdded}{326}
\newcommand{\BaseFalse}{2.30}
\newcommand{\BaseUnsupported}{3.4}
\newcommand{\BaseMisleading}{2.2}
\newcommand{\BaseBoardFalse}{0.90}
\newcommand{\CandFalse}{1.40}
\newcommand{\CandUnsupported}{5.1}
\newcommand{\CandMisleading}{1.9}
\newcommand{\CandBoardFalse}{0.90}
\newcommand{\PubFalse}{11.7}
\newcommand{\PubBoardFalse}{2.5}
\newcommand{\FalseDiff}{$-$0.90 {\footnotesize ($-$2.30 to +0.60)}}
\newcommand{\BoardDiff}{0.00 {\footnotesize ($-$1.00 to +1.00)}}
\newcommand{\FirstGradedPeriod}{week 35}
\newcommand{\FirstTruth}{preferred the candidate}
\newcommand{\NowTruth}{preferred the candidate}
\newcommand{\FirstCoaching}{preferred the candidate}
\newcommand{\NowCoaching}{preferred the candidate}
\newcommand{\FirstListening}{preferred today's prompt}
\newcommand{\NowListening}{preferred today's prompt}
\newcommand{\FirstOverall}{preferred the candidate}
\newcommand{\NowOverall}{preferred the candidate}
\newcommand{\RegradedN}{four}
\newcommand{\RegradeShare}{10}
\newcommand{\RegradeMax}{26}
\newcommand{\RuleTable}{%
\begin{tabular}{lrl}
\toprule
Condition & Estimate (95\% interval) & Holds \\
\midrule
Coaching on the train periods above 0 & +1.00 {\footnotesize (+1.00 to +1.00)} & yes \\
Coaching on the held-out periods above 0 & +1.00 {\footnotesize (+1.00 to +1.00)} & yes \\
Overall not below 0 & +0.60 {\footnotesize (+0.20 to +1.00)} & yes \\
False claims per script not up & $-$0.90 {\footnotesize ($-$2.30 to +0.60)} & yes \\
False board claims per script not up & 0.00 {\footnotesize ($-$1.00 to +1.00)} & yes \\
The check passes as often & 10 of 10 {\footnotesize (today's: 10)} & yes \\
Controls clean (noise floor: coaching) & 0.00 {\footnotesize ($-$0.75 to +0.75)} & yes \\
Coaching moves made more often & +0.36 {\footnotesize (+0.24 to +0.47)} & yes \\
At most 1.3 times as long & 1.37 {\footnotesize (1.20 to 1.52)} & \textbf{no} \\
\bottomrule
\end{tabular}
}
\newcommand{\Kept}{no}
\newcommand{\PeriodTable}{%
\begin{tabular}{llrccccrrrr}
\toprule
 & & & \multicolumn{4}{c}{Judge, from the candidate's side} & \multicolumn{2}{c}{Minutes heard once} & \multicolumn{2}{c}{Coaching share} \\
\cmidrule(lr){4-7}\cmidrule(lr){8-9}\cmidrule(lr){10-11}
Period & Split & Games & Truth & Coaching & Listening & Overall & Today & Candidate & Today & Candidate \\
\midrule
Week 35 & train & 55 & + & + & $-$ & + & 6.5 & 9.5 & 40 & 73 \\
Week 37 & train & 19 & 0 & + & $-$ & + & 6.1 & 9.3 & 50 & 100 \\
Week 39 & train & 59 & $-$ & + & 0 & $-$ & 6.2 & 7.7 & 50 & 69 \\
Week 41 & train & 34 & 0 & + & 0 & + & 6.5 & 10.4 & 60 & 69 \\
August & train & 307 & + & + & + & + & 7.3 & 5.6 & 44 & 56 \\
Week 36 & held out & 32 & 0 & + & 0 & 0 & 5.6 & 8.2 & 40 & 59 \\
Week 38 & held out & 49 & $-$ & + & $-$ & 0 & 6.1 & 9.3 & 40 & 93 \\
Week 40 & held out & 33 & + & + & + & + & 6.7 & 8.3 & 40 & 100 \\
Week 42 & held out & 34 & $-$ & + & + & + & 7.0 & 8.5 & 40 & 100 \\
September & held out & 172 & 0 & + & $-$ & + & 5.6 & 9.4 & 30 & 71 \\
\bottomrule
\end{tabular}
}


\title{Do the weekly and monthly videos coach better when their prompts ask for a coach's moves?\\[0.3em]
\large The daily video's bullet, adapted and added to the weekly and monthly routines' prompts,
replayed on \Cases{} past periods}
\author{Chess Coach research lab}
\date{}

\begin{document}
\maketitle

\begin{abstract}
The app's weekly and monthly routines write a narrated video beside each period's review. The
daily video's prompt took five moves of a strong human coach after a blind test: ask his move first,
name the habit behind the mistake and the habit that replaces it, play the better move on to the
opponent's best reply, say whether the miss is normal at his level, and praise a good move by what it
saw. We added the same bullet, adapted to the period video, to the weekly and monthly prompts and
replayed their review and video step on \Cases{} past periods (\WeeksN{} weeks and \MonthsN{}
months), once under each prompt, with the routines' own model and checker. Judging blind in both
orders, a model preferred the candidate's coaching on all \Cases{} periods, train and held out alike,
and its videos overall on \OverallWins{} periods against \OverallLosses. An audit against the evidence
found the candidate asking his move first in \CandOpenQuestion\% of its moment scenes (today's
prompts: \BaseOpenQuestion\%), playing the better move on to their reply in \CandThenWhat\%
(\BaseThenWhat\%) and placing the miss at his level in \CandToldLevel\% of the positions it teaches
(\BaseToldLevel\%), with no more false claims, and the judge's listening verdicts were even.
Every condition of the rule fixed before the first run held but one: heard once, the candidate's
videos are \LengthRatio{} times as long as today's, past the bound of \LengthCap, so the rule keeps
today's prompts. In the app's own voice the candidate's weekly videos last \CandWeekRendered{}
minutes, inside the five to eight the video grammar asks for, where today's last
\BaseWeekRendered. On the owner's call, both prompts now carry the bullet as tested (chess
PR~\PortPR).
\end{abstract}

\section{The question}

The app's weekly and monthly routines have a model review the player's games of the period and
write, beside the written review, a narrated video script, which the app renders as boards, charts,
captions and a synthetic voice. A period video synthesizes: its grammar asks for two or three board
scenes that show the period's recurring pattern through its clearest example, with how often it came
back, among chart, tile and bullet scenes, and closes on the period's one focus.

The lab's daily study (\emph{video-coaching}) tested five moves of a strong human coach in the daily
video: ask before telling; name the habit behind the mistake, then the habit that replaces it; follow
the better move to the opponent's best reply; say whether a miss is normal at the player's level; and
praise a good move by what it saw. The daily prompt took them. The weekly and monthly prompts say how
the video binds its boards and numbers but ask for none of the five. The candidate adds the daily's
bullet to step~4 of both prompts, the step that writes the annotations and the video script, with
three changes: it covers every board scene that teaches a move of his; the habit sub-point also says
how often the habit came back in the period; and the engine's reply and the findability come from the
day's archived analysis, since a period has no day's digest:

\CandidateBullet

Nothing else in the prompts changes. The question is whether the period video coaches better with
it, without becoming less true or much longer.

\section{The cases and the writer}

The cases are the \Cases{} periods with a review and a video script on the app's main branch from
its command-line era: \CaseSpan. Each is a checkout of the app at \texttt{\AppCommit}, with the
coach data as main held it just before that period's review landed, and the period's brief prepared
again by the app's own \texttt{coach weekly prepare} or \texttt{coach monthly prepare}, offline. A
week holds \WeekGamesMin{} to \WeekGamesMax{} games and the months \MonthGames. The split was fixed
before any run, interleaved by date: train is \TrainCases; held out is \HeldCases.

The daily study's writer wrote the video alone, beside the review the day published. That does not
work for a period: \RefFailedN{} of the \AllPeriods{} published reviews fail today's strict check
against their prepared brief (\RefErrorsMin{} to \RefErrorsMax{} errors), mostly for numbers the
brief prepared again no longer carries and, in the older reviews, for sections the check now
requires; only \RefPassed's passes. So each run writes the whole of step~4, the review and its video script
together, as the live routine does, and the graders read only the video script.

The writer runs as the routines do: Opus 5.5, headless, with nothing inherited from the session that
launched it, and a short preamble that names the period, says the brief is ready and stops the writer
once the strict check passes (no render, publish or commit). Each prompt wrote every period once, and
today's prompts wrote \NullCases{} a second time, to measure the judge's noise floor (\WriterRuns{}
runs in all).

\section{The graders}

A script with no model reads each run: whether the final check passes, the check rounds, commands it
should not have run, the script's shape, its length heard once (spoken words at \WordsPerSecond{} a
second plus every pause), cost and minutes.

Sonnet 5.5 audits each script alone against the evidence it was written from: the brief without its
bulk, the period's turning points with his think times, a line per game, the period's and the
career's digests, the week's daily logs or the month's weekly reviews, and for every board scene the
position as the archive replays it from the day the board names, what was played there and next, the
engine's best move and reply, the findability and the engine lines the day archived at that turning
point, how players at his band play there (Maia-2), the plan lines, the threats, what each piece hits
and guards, what each move the scene plays does, after each move the game did not play the engine's
best reply, and the time he took on each of his moves in the scene's game with the clock he had left.
It lists the claims the evidence contradicts (\emph{false}), does not carry (\emph{unsupported}) or
carries but slants (\emph{misleading}), and for each board scene its kind (a moment, praise, a
judgement call, a plan, other), its false claims about the board, and which of the five moves it
makes.

Sonnet 5.5 also judges each pair blind, in both orders, on four criteria: truth, coaching (with the
period's own duty first: name what repeated and show it on the board), listening (would he follow it
once, at speaking pace) and overall. A side wins a criterion only when both orders agree. Two
controls check the judge on every period: today's script against itself, and against a decoy that
adds as many spoken words as the candidate added on that period (at least \MinDecoyWords), in beats
that restate their scene with no square, move, piece, number or question the scene does not already
have.

The pilot wrote \FirstGradedPeriod{} under both prompts. Its candidate cited the long thinks that led
to a turning point, which the evidence then could not check, so the evidence gained his time on every
move of a scene's game, and every period, \FirstGradedPeriod{} again included, was graded with it.
The regrade's verdicts are the ones reported (Section~\ref{sec:twice}).

\section{The rule}

A period scores $+1$ when the candidate wins a criterion, $-1$ when it loses and $0$ for a tie;
intervals are 95\% percentile bootstrap intervals over periods. The candidate replaces the live
prompts only when every condition of Table~\ref{tab:rule} holds, fixed before its first run.

\begin{table}[h]
\centering
\small\RuleTable
\caption{The rule, fixed before the candidate's first run. Coaching and overall are the judge's
scores per period; the claim rows and the coaching moves are the audit's numbers per script, candidate
minus today's prompts; the length is the candidate's script heard once over today's on the same
period.}
\label{tab:rule}
\end{table}

\section{What the bullet changed}

\begin{table}[h]
\centering
\small\MovesTable
\caption{The coaching moves the audit finds, pooled over the \Cases{} periods: the share of the
scenes that owe a move and make it, with the number owing it in brackets. A moment scene owes the
habit, its replacement and the level; it owes the open question only when it asks him something, and
the reply only when it shows a better move. In the middle rows, consecutive board scenes on the same
position count once, and a move counts as made when any of them makes it.}
\label{tab:moves}
\end{table}

Today's prompts name the habit behind a mistake and end on its replacement in most moment scenes,
as the published scripts do, but they rarely ask his move before giving a hint
(\BaseOpenQuestion\%), rarely play the better move on to the opponent's best reply
(\BaseThenWhat\%) and rarely say whether a miss is normal at his level (\BaseLevel\%). The candidate
does all three in most moment scenes (\CandOpenQuestion\%, \CandThenWhat\% and \CandLevel\%), the
level it states agrees with the findability evidence in \CandLevelTrue\% of the scenes that state
one, and every praise scene it writes names what the move saw. Over the five moves, the share the
audit finds in the moment scenes rose from \ShareBase\% to \ShareCand\% (\ShareDiff{} a script;
rule~7). It also shows more boards (\CandBoardScenes{} board scenes a script against
\BaseBoardScenes) and teaches more moments in them (\CandMoments{} against \BaseMoments).

Scene by scene, the candidate names the habit and its replacement a little less often
(\CandHabitNamed\% and \CandReplacement\%), for the reason the daily study found: it often shows a
position in two scenes, one that asks and one that answers (\CandSplit{} board scenes a script repeat
the position of the board scene before them, against \BaseSplit{} under today's prompts), and the
scene that asks counts as a miss. Counted once per position taught, it names the habit and ends on
its replacement in \CandToldHabitNamed\% and \CandToldReplacement\% of them, against
\BaseToldHabitNamed\% and \BaseToldReplacement\%.

\section{What the judge preferred}

\begin{table}[h]
\centering
\small
\begin{tabular}{lrrr}
\toprule
Criterion & Train (\TrainN{} periods) & Held out (\HeldN{} periods) & All \\
\midrule
Truth & \TrainTruth & \HeldTruth & \AllTruth \\
Coaching & \TrainCoaching & \HeldCoaching & \AllCoaching \\
Listening & \TrainListening & \HeldListening & \AllListening \\
Overall & \TrainOverall & \HeldOverall & \AllOverall \\
\bottomrule
\end{tabular}
\caption{The judge's mean score per period from the candidate's side ($+1$ a win, $-1$ a loss, $0$ a
tie), with its 95\% bootstrap interval over periods.}
\label{tab:judge}
\end{table}

The judge preferred the candidate's coaching on all \Cases{} periods, in both orders each time, and
it preferred the candidate overall on \OverallWins{} periods against \OverallLosses, with
\OverallTies{} ties (Table~\ref{tab:judge}). Truth was even: \TruthWins{} wins, \TruthLosses{}
losses and \TruthTies{} ties. So was listening: \ListeningWins{} wins, \ListeningLosses{} losses
and \ListeningTies{} ties. The daily study's candidate never won listening; this one won it about
as often as it lost it.

Two runs of today's prompts on the same week, judged against each other, split on coaching on
\NullDecisive{} of \NullN{} weeks, once each way, for a mean of \AllNullCoaching. The judge does pick
a side between two scripts of one prompt, but not one way on every period, as it did for the
candidate. Read in one order only, the judge favoured the side it read second in \SecondShare\% of
its \Decided{} decided verdicts, so a verdict that only one order gives is a tie.

\section{The controls}

Today's script judged against itself tied on every criterion on all \SameN{} periods. The decoy,
today's script with on average \DecoyAdded{} spoken words added that only restate their scene, never
won coaching (it lost it on \DecoyLostCoaching{} periods) and lost listening and overall on all
\DecoyN. The judge does not reward words that say nothing new, so the candidate's wins come from what
its extra words say.

\section{Truth}

The audit found fewer false claims in the candidate's scripts than in today's (\CandFalse{} and
\BaseFalse{} a script; the difference is \FalseDiff), within the noise, and as many false claims
about the board (\CandBoardFalse{} against \BaseBoardFalse{} a script; \BoardDiff), though the
candidate shows more board scenes. It carries more lines the evidence does not settle
(\CandUnsupported{} unsupported a script against \BaseUnsupported) and about as many misleading ones
(\CandMisleading{} against \BaseMisleading).

\section{Length, cost and the check}

\begin{table}[h]
\centering
\small\LengthTable
\caption{Minutes per video: heard once (spoken words at \WordsPerSecond{} a second plus every pause,
the rule's measure) and rendered (the app's plan voiced by its own voice, the measure the video
grammar's minutes use), and how many rendered videos fall inside the grammar's minutes.}
\label{tab:length}
\end{table}

Every run of both prompts passed the routine's check. Heard once, the candidate's videos last \CandHeard{} minutes against \BaseHeard{}, with
\CandWords{} spoken words against \BaseWords{} and \CandQuestions{} questions against
\BaseQuestions; on the same period it is \LengthRatioInterval{} times as long, past the
\LengthCap{} rule~8 allows. Only on August was the candidate's video the shorter.

The rule measures length heard once, at a narrator's pace. The video grammar gives its minutes for
the rendered video, and the app's voice is quicker: rendered, a video lasts on average
\RenderedShare\% of its length heard once (the app's own estimate of it, from the words alone, is
within \EstimateOff\% of the voice on every script). Rendered, the candidate's weekly videos last
\CandWeekRendered{} minutes and today's \BaseWeekRendered{} (Table~\ref{tab:length}); on the same
period the candidate is \RenderedRatio{} times as long rendered, about as much as heard once. The
grammar asks for \WeekTargetLo{} to \WeekTargetHi{} minutes a week: \CandWeekRenderedIn{} of the
candidate's eight weeks fall inside it and \CandWeekRenderedOver{} runs past it, against
\BaseWeekRenderedIn{} of today's, whose other \BaseWeekRenderedUnder{} run short. The months are two
and say little: \MonthLengthText. This comparison was made after the result and decides nothing.

The writer took \CandRounds{} check rounds against \BaseRounds, \CandMinutes{} minutes against
\BaseMinutes{} and \$\CandCost{} against \$\BaseCost{} a period at list price.

\section{The graders measured twice}
\label{sec:twice}

The pilot's \RegradedN{} scripts (week~35 under both prompts and the published scripts of week~35
and August) were audited with the first evidence and again with the final one. Their coaching share
moved by \RegradeShare{} points on average and by up to \RegradeMax, so one period's share in
Table~\ref{tab:periods} says little, while the pooled difference (\ShareDiff) is well beyond that.
The judge's verdicts on week~35 were the same on every criterion both times.

\begin{table}[h]
\centering
\scriptsize\PeriodTable
\caption{Each period: the judge's verdict from the candidate's side ($+$ a win, $-$ a loss, $0$ a
tie), the minutes heard once under each prompt, and the share of the coaching moves the audit finds
in the moment scenes.}
\label{tab:periods}
\end{table}

\section{What the app takes}

The rule keeps today's prompts: the candidate held every condition but the length
(Table~\ref{tab:rule}). It won the coaching criterion on every period, with no more false claims and
even listening verdicts, and its extra length comes with the three moves today's prompts skip: a
question before each answer, the reply played out, and the level. Whether that length is worth it was
the owner's call, which the rule, fixed before the first run, does not make. The owner took it: the
weekly and monthly prompts carry the bullet as tested (chess PR~\PortPR). A next round could ask for
the same moves in fewer words, as the daily study's next round would.

\section{Limits}

\begin{itemize}[nosep]
\item The five moves are one change: the result shows they help together, not which of them does.
\item Each run writes the period's review and its video together, so two runs of a period may show
different boards; the judge compares two videos, not two treatments of one board.
\item A replay writes the review and the video after the brief is in place; the routine writes them
at the end of a longer run. Over the last four weeks (\RecentWeeks) the published videos ran about as
long heard once as today's replays of the same weeks (\PubRecentHeard{} against \BaseRecentHeard{}
minutes), so the replays' lengths look like the routine's; the earlier published weeks ran longer.
\item The graders are models. The judge picks a side between two runs of one prompt on some periods
and favours the side it reads second; the audit's tallies move between two readings. The controls,
the noise floor and both orders bound these, and the graders are a different model from the writer.
\item Ten periods, two of them months: the months' results stand on one period each.
\item The level a script states is checked against the same findability the writer reads (Maia-2 at
his band), so the audit checks that the writer read it right, not that Maia-2 knows his level.
\item The eval spent \$\Spent{} at list price (\$\WriterSpend{} on the writer, \$\GraderSpend{} on
the graders, \$\LostSpend{} on runs a container restart cut short), within its cap of \$\Cap, on the
owner's subscription, whose seven-day window it moved from \SevenDayFirst\% to \SevenDayLast\%.
\end{itemize}


\end{document}
